% =============================================================================
%  §II.2 — Couples et produit cartésien
%  Fragment inclus par main.tex (\input). Blocs subsection seulement.
% =============================================================================

\subsection{§II.2 --- Caractérisation du couple}
\textbf{Énoncé (Bourbaki, E~II.7).} $z=(x,y)\Leftrightarrow\bigl(z\text{ est un couple}
\ \text{et}\ x=\mathrm{pr}_1 z\ \text{et}\ y=\mathrm{pr}_2 z\bigr)$.

\textbf{Formalisé (CLOS).} L'équivalence, « $z$ couple » $:=(\exists a)(\exists b)(z=(a,b))$
inline. Le piège de capture (les lettres $x,y$ de l'énoncé coïncident avec les liants
$\tau$ de $\mathrm{pr}_1,\mathrm{pr}_2$) est levé par des \emph{témoins canoniques}
(variables libres $\subseteq\{z\}$). Module : \texttt{.../ii\_2\_couples\_produit/ensembles\_couple\_caracterisation.py}.

\subsection{§II.2 --- $z=(\mathrm{pr}_1 z,\mathrm{pr}_2 z)\Leftrightarrow z$ couple}
\textbf{Énoncé (Bourbaki, E~II.7).} « La relation $z=(\mathrm{pr}_1(z),\mathrm{pr}_2(z))$
est équivalente à "$z$ est un couple". »

\textbf{Formalisé (CLOS).} Corollaire de la caractérisation précédente, instanciée en
$x:=\mathrm{pr}_1 z$ et $y:=\mathrm{pr}_2 z$ : le membre droit $\bigl(z\text{ couple}\wedge
\mathrm{pr}_1 z=\mathrm{pr}_1 z\wedge\mathrm{pr}_2 z=\mathrm{pr}_2 z\bigr)$ se réduit à
« $z$ couple » par réflexivité des deux égalités. \texttt{couple\_egal\_projections}.
